% Section 3.1 -- what the extractor is, compressed. Mechanics (terminator set,
% tokenizer, namespace prefixes, keyword list) live in Appendix A.
% No literal numbers: every figure is a macro from numbers.tex.

\subsection{Source-visible extraction at declaration scope}
\label{sec:extraction}

The extractor takes two inputs: the mathlib checkout at the pinned revision,
and the set of declaration names in the published graph. It compiles nothing
and elaborates nothing, which is what makes it an instrument for the explicit
subgraph rather than an approximation of the full one.

Its scope is the whole declaration, statement and proof. Explicit edges come
from the elaborated term of the entire declaration, so a hypothesis contributes
premises. In \texttt{theorem foo (h : Nat.Prime p)}, \texttt{Nat.Prime} is
an explicit target. Extracting only what follows \texttt{:=} would drop
such targets while leaving them in the denominator, which is an under-scoped
instrument rather than a conservative one. The release ships no source
positions, so the block is located textually; the locate rate is measured
rather than assumed, at \auditDeclsLocated{} of \auditDeclsDrawn{} sampled
declarations.

Within that block the identifier tokens are resolved against the declaration
set, under the block-location, tokenization, resolution and prefix rules
Appendix~\ref{app:extractor} gives in full.

Two failure modes matter for reading the results. The first is deliberate: the
extractor never searches the declaration set by suffix, because resolving
\texttt{mp} that way would credit every declaration in mathlib whose name ends
in \texttt{.mp}. Tokens that resolve to nothing therefore stay out of the
premise set entirely, and the cost of that refusal is the largest component of the
shortfall below the declaration reference point of Section~\ref{sec:retention}
(\residualResolutionPct{}\%, Appendix~\ref{app:residual}). The second is
imposed: dot notation names a declaration only once the type of the local it
projects from is known, which requires elaboration, so that residue is reported
rather than resolved.
